\section{Reviewer A}

Reviewer comments:
{\it 
``The manuscript presents a new method to improve angular uncertainty in antineutrino direction reconstruction for segmented detectors. The authors apply a previously published algorithm to simulated data, optimize it for segmented detector geometries, and ultimately compare the results with observed data. Compared to conventional methods, the proposed approach aims to optimize for low-count limitations and provides a more robust way to handle error estimation.

The technical details are sound, and the results are clearly presented across various conditions, such as different segment sizes. The authors also provide a helpful elaboration of the assumptions underlying this work at the end of page 5. However, the manuscript would benefit from further clarification regarding the potential impact of these factors on the angular uncertainty estimation and the overall reconstruction. Given that these assumptions may have non-trivial effects, providing additional detail would significantly help to strengthen the robustness of this work.
''
}

Thank you for your comments and helpful feedback. We appreciate the positive assessment and your suggestions for improving the work. To that end, our response to your comments addresses the assumptions on page 5 listed below:

\begin{itemize}
    \item impact of 1/\(r^2\) assumption
    \item impact of angular spread of source
    \item impact of escaping neutrons \\ Neutrons that escape the detector are not considered because they cannot be binned. But also, neutrons that escape the binning radius (i.e. more than 1 voxel away from the prompt event in 150mm segment size, more than 2 voxels away from the prompt event in 50mm segment size, and more than 50 voxels away in the 5mm segment size), even if detected, will not be considered for directional reconstruction. This exclusion of events for the 5mm segment size had the greatest impact on efficiency. Even so, approximately 95\% of IBD events were binned. This is more of a concern for a real detector, where IBD near the surface is more likely to result in a neutron escaping. % wrap this up with how it impacts uncertainty (it's already accounted for)
\end{itemize}

The following statement has been added to the conclusion highlighting a planned use case study in which the impacts of these assumptions are relevant and will be studied:

``Applying this method to different deployment scenarios, such as small modular reactors, will investigate some of the underlying assumptions, such as neglecting \(\frac{1}{r^2}\) dependence and non-zero angular size of nearby reactors are all being considered for future work.''

Sincerely,
Brian Crow, University of Hawai'i at Manoa% fix the formatting for affiliation
on behalf of the authors

\section{Reviewer B}

Thank you for your comments and we appreciate your feedback. Rest assured that we understand that the end of semester can be a busy time for reviewers, so we appreciate you getting back to us. We also appreciate the recommendation for publication in PR Applied, and will incorporate your suggestions in our revised manuscript. Please find our responses to your comments in the order presented:

\begin{itemize}
    \item Abstract \\ The abstract has been changed to read: ``We treat data from Lithium-6-doped detector segments in the form of a matrix.'' -- specifying that the simulated detector data are from a lithium-6-doped detector. Our intent to generalize the method to other detector implementations is addressed in the conclusion.
    \item Page 2 \\ The following statement has been added explaining that the limited spatial distribution of the neutron detection layers in SOLID and CHANDLER was not studied during the development of this algorithm: \\
    \begin{quote}
         This confinement of the ${}^6$Li dopant to the surfaces of the segments in these detectors makes the neutron detection layer much thinner. This thin neutron detection layer contradicts the implicit assumption of isotropy in the binning scheme introduced in Figs. \ref{fig_segmentation_scaling} and \ref{fig:phi_def}. Modifying the binning and the directionality algorithm to address these anisotropies is not within the scope of this work.
    \end{quote}
    \item Conclusion \\ The planned studies to apply the method to other detector technologies and search for ways to improve mathematical robustness has been more explicitly stated with the following final sentence: \\
    \begin{quote}
        Planned future studies will investigate incorporating various statistical methods in directional analysis, applying this method to detectors of different geometry such as monolithic and separated, as well as with different detection media such as water Cherenkov, and applying the method to three dimensionally segmented detectors. 
    \end{quote}
\end{itemize}